% Part: normal-modal-logic
% Chapter: tableaux
% Section: completeness

\documentclass[../../../include/open-logic-section]{subfiles}

\begin{document}

\olfileid[hi]{nml}{tab}{cpl}

\olsection{\Log{K} के लिए पूर्णता}

\begin{explain}
  !!{tableau} की विधि की पूर्णता दिखाने के लिए हमें दिखाना है कि जब भी
  $\Gamma \Proves !A$ दिखाने वाला कोई संवृत !!{tableau} न हो, तब
  $\Gamma \Entails/ !A$, अर्थात् कोई प्रतिप्रतिरूप हो। किंतु ``कोई संवृत
  !!{tableau} नहीं है'' का अर्थ है कि उसे बनाने का हर संभावित प्रयास संवृत
  होने में विफल होना चाहिए। युक्ति यह देखना है कि यदि ऐसा हर ढंग संवृत होने
  में विफल हो, तो कोई विशिष्ट, \emph{व्यवस्थित और समग्र} ढंग भी विफल होगा।
  और यदि कोई संवृत !!{tableau} विद्यमान हो, तो यह व्यवस्थित और समग्र ढंग
  संवृत हो जाएगा। एक ही टैब्लो की खुली शाखाओं में प्रतिप्रतिरूप परिभाषित
  करने के लिए आवश्यक सारी सूचना होगी। इस टैब्लो की किसी खुली शाखा से मिला
  प्रतिप्रतिरूप उस शाखा पर प्रयुक्त सभी उपसर्गों को जगत बनाएगा, और कोई
  !!a{propositional variable}~$p$, $\sigma$ पर तभी सत्य होगा जब
  $\sFmla{\True}{p}[\sigma]$ शाखा पर आता हो।
\end{explain}

\begin{defn}
  किसी !!a{tableau} की शाखा \emph{पूर्ण} कहलाती है यदि जब भी उसमें ऐसा
  उपसर्गयुक्त !!{formula} $\sFmla{S}{!A}[\sigma]$ हो जिस पर कोई नियम लगाया
  जा सकता है, तब उसमें यह भी हो:
  \begin{enumerate}
    \item प्रतिज्ञप्तिक स्तंभन-नियमों की स्थिति में नियम के संबद्ध
      निष्कर्ष बने उपसर्गयुक्त !!{formula};
    \item प्रतिज्ञप्तिक शाखन-नियमों की स्थिति में संबद्ध निष्कर्ष
      !!{formula}ों में से कोई एक;
    \item नया उपसर्ग अपेक्षित करने वाले मोडल नियमों की स्थिति में
      कम-से-कम एक संभावित निष्कर्ष;
    \item प्रयुक्त उपसर्ग अपेक्षित करने वाले मोडल नियमों की स्थिति में
      शाखा पर आने वाले प्रत्येक उपसर्ग के लिए संबद्ध निष्कर्ष।
  \end{enumerate}
\end{defn}

\begin{explain}
उदाहरणार्थ, किसी पूर्ण शाखा में जब भी $\sFmla{\True}{!B \land !C}$ हो,
तब $\sFmla{\True}{!B}[\sigma]$ और $\sFmla{\True}{!C}[\sigma]$, दोनों
होते हैं। यदि उसमें $\sFmla{\True}{!B \lor !C}[\sigma]$ हो, तो उसमें
$\sFmla{\False}{!B}[\sigma]$ और $\sFmla{\True}{!C}[\sigma]$ में से
कम-से-कम एक होता है। यदि उसमें \iftag{prvBox}
{$\sFmla{\False}{\Box}[\sigma]$} {$\sFmla{\True}{\Diamond}[\sigma]$} हो,
तो उसमें कम-से-कम एक~$n$ के लिए \iftag{prvBox}
{$\sFmla{\False}{\Box}[\sigma.n]$}
{$\sFmla{\True}{\Diamond}[\sigma.n]$} भी होता है। और जब भी उसमें
\iftag{prvBox} {$\sFmla{\True}{\Box}[\sigma]$}
{$\sFmla{\False}{\Diamond}[\sigma]$} हो, तब शाखा पर $\sigma.n$ प्रयुक्त
होने वाले प्रत्येक~$n$ के लिए उसमें \iftag{prvBox}
{$\sFmla{\True}{\Box}[\sigma.n]$}
{$\sFmla{\False}{\Diamond}[\sigma.n]$} भी होता है।
\end{explain}

\begin{prop}\ollabel{prop:complete-tableau}
  प्रत्येक परिमित $\Gamma$ का ऐसा !!a{tableau} है जिसमें हर शाखा पूर्ण है।
\end{prop}

\begin{proof}
  $\Gamma$ के किसी !!a{tableau} की खुली शाखा पर विचार करें। शाखा में ऐसे
  परिमिततः अनेक उपसर्गयुक्त !!{formula} हैं जिन पर कोई नियम लगाया जा सकता
  है। किसी नियत क्रम (जैसे ऊपर से नीचे) में, इन उपसर्गयुक्त !!{formula}ों में
  से प्रत्येक ऐसे सूत्र के लिए जिस पर शर्तें (1)--(4) पहले से सत्य नहीं
  हैं, शाखा बढ़ाने हेतु उस पर लागू हो सकने वाले नियम लगाएँ। कुछ स्थितियों
  में इससे शाखन होगा; बच गए सभी उपसर्गयुक्त !!{formula}ों के लिए प्रत्येक
  परिणामी शाखा के सिरे पर नियम लगाएँ। चूँकि उपसर्गयुक्त !!{formula}ों की
  संख्या परिमित है और शाखा पर प्रयुक्त उपसर्गों की संख्या परिमित है, यह
  प्रक्रिया अंततः मूल शाखा को विस्तृत करने वाली (संभवतः अनेक) शाखाएँ देती
  है। प्रत्येक पर यह प्रक्रिया लगाएँ और दोहराएँ। किंतु निर्माण से हर शाखा
  संवृत है।
\end{proof}

\begin{thm}[पूर्णता]
  \ollabel{thm:tableau-completeness}
  यदि $\Gamma$ का कोई संवृत !!{tableau} नहीं है, तो $\Gamma$ संतोषणीय है।
\end{thm}

\begin{proof}
प्रतिज्ञप्ति से $\Gamma$ का ऐसा !!a{tableau} है जिसमें हर शाखा पूर्ण है।
चूँकि उसका कोई संवृत !!{tableau} नहीं है, उसका ऐसा !!a{tableau} है जिसमें
कम-से-कम एक शाखा खुली और पूर्ण है। $\Delta$ उस शाखा के उपसर्गयुक्त
!!{formula}ों का समुच्चय हो और $P(\Delta)$ उसमें आने वाले उपसर्गों का
समुच्चय हो।

हम प्रतिरूप~$\mModel{M(\Delta)} = \tuple{P(\Delta), R, V}$ परिभाषित
करते हैं, जहाँ जगत $\Delta$ में आने वाले उपसर्ग हैं और अभिगम्यता संबंध
यह है:
\[
R\sigma\sigma' \quad \text{तभी जब} \quad
\sigma'=\sigma.n \quad \text{किसी~$n$ के लिए}
\]
और
\[
V(p) = \Setabs{\sigma}{\sFmla{\True}{p}[\sigma] \in \Delta}.
\]
हम $!A$ पर आगमन से दिखाते हैं कि यदि $\sFmla{\True}{!A}[\sigma] \in
\Delta$, तो $\mSat{M(\Delta)}{!A}[\sigma]$; और यदि
$\sFmla{\False}{!A}[\sigma] \in \Delta$, तो
$\mSat/{M(\Delta)}{!A}[\sigma]$।

\begin{enumerate}
  \item \indcase{!A}{p}{यदि $\sFmla{\True}{\indfrm}[\sigma] \in
    \Delta$, तो $V$ की परिभाषा से $\sigma \in V(p)$, और इसलिए
    $\mSat{M(\Delta)}{\indfrm}[\sigma]$।

    यदि $\sFmla{\False}{\indfrm}[\sigma] \in \Delta$, तो
    $\sFmla{\True}{\indfrm}[\sigma] \notin \Delta$, क्योंकि अन्यथा शाखा
    संवृत होती। अतः $\sigma \notin V(p)$ और इसलिए
    $\mSat/{M(\Delta)}{\indfrm}[\sigma]$।}
  \item \indcase{!A}{\lnot !B}{\iftag{probNot}{अभ्यास।}{यदि
      $\sFmla{\True}{\indfrm}[\sigma] \in \Delta$, तो शाखा के पूर्ण होने
      से $\sFmla{\False}{!B}[\sigma] \in \Delta$। आगमनात्मक परिकल्पना से
      $\mSat/{M(\Delta)}{!B}[\sigma]$, अतः
      $\mSat{M(\Delta)}{\indfrm}[\sigma]$।

      यदि $\sFmla{\False}{\indfrm}[\sigma] \in \Delta$, तो शाखा के पूर्ण
      होने से $\sFmla{\True}{!B}[\sigma] \in \Delta$। आगमनात्मक परिकल्पना
      से $\mSat{M(\Delta)}{!B}[\sigma]$, अतः
      $\mSat/{M(\Delta)}{\indfrm}[\sigma]$।}}
  \item \indcase{!A}{!B \land !C}{\iftag{probAnd}{अभ्यास।}{यदि
      $\sFmla{\True}{\indfrm}[\sigma] \in \Delta$, तो शाखा के पूर्ण होने
      से $\sFmla{\True}{!B}[\sigma] \in \Delta$ और
      $\sFmla{\True}{!C}[\sigma] \in \Delta$, दोनों। आगमनात्मक परिकल्पना
      से $\mSat{M(\Delta)}{!B}[\sigma]$ और
      $\mSat{M(\Delta)}{!C}[\sigma]$। अतः
      $\mSat{M(\Delta)}{\indfrm}[\sigma]$।

      यदि $\sFmla{\False}{\indfrm}[\sigma] \in \Delta$, तो शाखा के पूर्ण
      होने से या तो $\sFmla{\False}{!B}[\sigma] \in \Delta$, या
      $\sFmla{\False}{!C}[\sigma] \in \Delta$। आगमनात्मक परिकल्पना से या
      तो $\mSat/{M(\Delta)}{!B}[\sigma]$, या
      $\mSat/{M(\Delta)}{!B}[\sigma]$। अतः
      $\mSat/{M(\Delta)}{\indfrm}[\sigma]$।}}
  \item \indcase{!A}{!B \lor !C}{\iftag{probOr}{अभ्यास।}{यदि
      $\sFmla{\True}{\indfrm}[\sigma] \in \Delta$, तो शाखा के पूर्ण होने
      से या तो $\sFmla{\True}{!B}[\sigma] \in \Delta$, या
      $\sFmla{\True}{!C}[\sigma] \in \Delta$। आगमनात्मक परिकल्पना से या
      तो $\mSat{M(\Delta)}{!B}[\sigma]$, या
      $\mSat{M(\Delta)}{!C}[\sigma]$। अतः
      $\mSat{M(\Delta)}{\indfrm}[\sigma]$।

      यदि $\sFmla{\False}{\indfrm}[\sigma] \in \Delta$, तो शाखा के पूर्ण
      होने से $\sFmla{\False}{!B}[\sigma] \in \Delta$ और
      $\sFmla{\False}{!C}[\sigma] \in \Delta$, दोनों। आगमनात्मक परिकल्पना
      से $\mSat/{M(\Delta)}{!B}[\sigma]$ और
      $\mSat/{M(\Delta)}{!B}[\sigma]$, दोनों। अतः
      $\mSat/{M(\Delta)}{\indfrm}[\sigma]$।}}
  \item \indcase{!A}{!B \lif !C}{\iftag{probIf}{अभ्यास।}{यदि
      $\sFmla{\True}{\indfrm}[\sigma] \in \Delta$, तो शाखा के पूर्ण होने
      से या तो $\sFmla{\False}{!B}[\sigma] \in \Delta$, या
      $\sFmla{\True}{!C}[\sigma] \in \Delta$। आगमनात्मक परिकल्पना से या
      तो $\mSat/{M(\Delta)}{!B}[\sigma]$, या
      $\mSat{M(\Delta)}{!C}[\sigma]$। अतः
      $\mSat{M(\Delta)}{\indfrm}[\sigma]$।

      यदि $\sFmla{\False}{\indfrm}[\sigma] \in \Delta$, तो शाखा के पूर्ण
      होने से $\sFmla{\True}{!B}[\sigma] \in \Delta$ और
      $\sFmla{\False}{!C}[\sigma] \in \Delta$, दोनों। आगमनात्मक परिकल्पना
      से $\mSat{M(\Delta)}{!B}[\sigma]$ और
      $\mSat/{M(\Delta)}{!B}[\sigma]$, दोनों। अतः
      $\mSat/{M(\Delta)}{\indfrm}[\sigma]$।}}
  \item \indcase{!A}{\Box !B}{\iftag{probBox}{अभ्यास।}{यदि
      $\sFmla{\True}{\indfrm}[\sigma] \in \Delta$, तो शाखा के पूर्ण होने
      से शाखा पर प्रयुक्त प्रत्येक $\sigma.n$ के लिए
      $\sFmla{\True}{!B}[\sigma.n] \in \Delta$, अर्थात् प्रत्येक ऐसे
      $\sigma' \in P(\Delta)$ के लिए जिसके लिए $R\sigma\sigma'$।
      आगमनात्मक परिकल्पना से प्रत्येक ऐसे $\sigma'$ के लिए
      $\mSat{M(\Delta)}{!B}[\sigma']$ जिसके लिए $R\sigma\sigma'$। अतः
      $\mSat{M(\Delta)}{\indfrm}[\sigma]$।

      यदि $\sFmla{\False}{\indfrm}[\sigma] \in \Delta$, तो शाखा के पूर्ण
      होने से किसी $\sigma.n$ के लिए
      $\sFmla{\False}{!B}[\sigma.n] \in \Delta$। आगमनात्मक परिकल्पना से
      $\mSat/{M(\Delta)}{!B}[\sigma.n]$। चूँकि $R\sigma(\sigma.n)$, ऐसा
      कोई~$\sigma'$ है कि $\mSat/{M(\Delta)}{!B}[\sigma']$। अतः
      $\mSat/{M(\Delta)}{\indfrm}[\sigma]$।}}
  \item \indcase{!A}{\Diamond !B}{\iftag{probDiamond}{अभ्यास।}{यदि
      $\sFmla{\True}{\indfrm}[\sigma] \in \Delta$, तो शाखा के पूर्ण होने
      से किसी $\sigma.n$ के लिए $\sFmla{\True}{!B}[\sigma.n] \in
      \Delta$। आगमनात्मक परिकल्पना से
      $\mSat{M(\Delta)}{!B}[\sigma.n]$। चूँकि $R\sigma(\sigma.n)$, ऐसा
      कोई~$\sigma'$ है कि $\mSat{M(\Delta)}{!B}[\sigma']$। अतः
      $\mSat{M(\Delta)}{\indfrm}[\sigma]$।

      यदि $\sFmla{\False}{\indfrm}[\sigma] \in \Delta$, तो शाखा के पूर्ण
      होने से शाखा पर प्रयुक्त प्रत्येक $\sigma.n$ के लिए
      $\sFmla{\False}{!B}[\sigma.n] \in \Delta$, अर्थात् प्रत्येक ऐसे
      $\sigma' \in P(\Delta)$ के लिए जिसके लिए $R\sigma\sigma'$।
      आगमनात्मक परिकल्पना से प्रत्येक ऐसे $\sigma'$ के लिए
      $\mSat/{M(\Delta)}{!B}[\sigma']$ जिसके लिए $R\sigma\sigma'$। अतः
      $\mSat/{M(\Delta)}{\indfrm}[\sigma]$।}}
\end{enumerate}
चूँकि $\Gamma \subseteq \Delta$, इसलिए $\mSat{M(\Delta)}{\Gamma}$।
\end{proof}

\begin{prob}
\olref[nml][tab][cpl]{thm:tableau-completeness} का प्रमाण पूरा कीजिए।
\end{prob}

\begin{cor}
\ollabel{cor:entailment-completeness}
यदि $\Gamma \Entails !A$, तो $\Gamma \Proves !A$।
\end{cor}

\begin{cor}
\ollabel{cor:weak-completeness}
यदि $!A$ सभी प्रतिरूपों में सत्य है, तो $\Proves !A$।
\end{cor}

\end{document}
